%% %%%%%%%%%%%%%%%%%%%%%%%%%%%%%%%%%%%%%%%%%%%%%%%%%%%%%%%%%%%%%%%
%% AGFA-Book-Kapitel-2.tex -- Kapitel II
%% %%%%%%%%%%%%%%%%%%%%%%%%%%%%%%%%%%%%%%%%%%%%%%%%%%%%%%%%%%%%%%%
\chapter{Anwendungen}
\label{chap:anwendungen}

\section{Differentialgleichungen}
\label{sec:dgl}

Wir wenden \cref{thm:banach} aus \cref{subsec:kontraktion} an
(siehe auch \cref{sec:vollstaendig} und \cref{cor:eindeutig}).
Mit \eqref{eq:apriori} erhält man eine Fehlerschranke.

\begin{theorem}[Picard--Lindelöf]
\label{thm:picard}
Ist $f$ Lipschitz-stetig, so hat das Anfangswertproblem genau eine Lösung.
\end{theorem}

\begin{proof}
\Cref{thm:banach} auf den Integraloperator anwenden; die Metrik aus
\cref{def:metrik} ist dabei die Supremumsmetrik.
\end{proof}

\begin{exercise}
\label{ex:picard}
Man führe die Abschätzung in \cref{thm:picard} aus.
\end{exercise}

Verweise auf mehrere Ergebnisse: \cref{lem:cauchy,thm:banach,thm:picard}.
